% GENERATED FILE. Edit canonical lessons, then run npm run generate:books.
% canonical-source-hash: fnv1a64:753bef2c2cac0322
% canonical-lessons: MR-C165-cukne, MR-C165-sudharne, MR-C165-bighadne, MR-C165-niriksan-karne, MR-C165-bhag-ghene

\chapter{To miss, To improve, To get worse, To observe, To take part}
\label{ch:mr-to-miss-to-improve-to-get-worse-to-observe-to-take-part}
\hlchaptermodality{165}

\begin{chapteropening}
\textbf{By the end of this chapter:} \emph{I can say to miss, to improve, to get worse, to observe and to take part.}

Complete the last lesson of chapter 165: I can say to miss, to improve, to get worse, to observe and to take part.
\end{chapteropening}

\section[\texorpdfstring{\mr{चुकणे} (cukṇe)}{चुकणे (cukṇe)}]{\mr{चुकणे} (cukṇe) — to miss; to go wrong}

\label{lesson:MR-C165-cukne}



\begin{quote}
Before the new one: say the Marathi for to increase, then the Marathi for to reduce.
\end{quote}

\subsection*{You'll want to know: \mr{चुकणे}}
\textbf{\mr{चुकणे}} — \emph{cukṇe} — \textquotedblleft{}to miss; to go wrong\textquotedblright{}.

\begin{grammarlens}[title={What you've built}]
One more word for this chapter.
\end{grammarlens}

\subsection*{Your turn}
\begin{itemize}
\raggedright
  \item \emph{Say it:} \emph{cukṇe}
  \item \emph{Say it:} \emph{cukṇe}, once more
  \item \emph{Say it:} say \emph{cukṇe}
  \item \emph{Recall:} say the Marathi for to increase, then the Marathi for to reduce, then say \emph{cukṇe} again
\end{itemize}

\begin{tcolorbox}[breakable,colback=teal!4,colframe=teal!35!black,title={Before you move on}]
What does \mr{चुकणे} mean? (\textquotedblleft{}To miss; to go wrong\textquotedblright{}.) Say it once more.
\end{tcolorbox}
\section[\texorpdfstring{\mr{सुधारणे} (sudhārṇe)}{सुधारणे (sudhārṇe)}]{\mr{सुधारणे} (sudhārṇe) — to improve}

\label{lesson:MR-C165-sudharne}



\begin{quote}
Before the new one: say the Marathi for to reduce, then the Marathi for to miss; to go wrong.
\end{quote}

\subsection*{You'll want to know: \mr{सुधारणे}}
\textbf{\mr{सुधारणे}} — \emph{sudhārṇe} — \textquotedblleft{}to improve\textquotedblright{}.

\begin{grammarlens}[title={What you've built}]
One more word for this chapter.
\end{grammarlens}

\subsection*{Your turn}
\begin{itemize}
\raggedright
  \item \emph{Say it:} \emph{sudhārṇe}
  \item \emph{Say it:} \emph{sudhārṇe}, once more
  \item \emph{Say it:} say \emph{sudhārṇe}
  \item \emph{Recall:} say the Marathi for to reduce, then the Marathi for to miss; to go wrong, then say \emph{sudhārṇe} again
\end{itemize}

\begin{tcolorbox}[breakable,colback=teal!4,colframe=teal!35!black,title={Before you move on}]
What does \mr{सुधारणे} mean? (\textquotedblleft{}To improve\textquotedblright{}.) Say it once more.
\end{tcolorbox}
\section[\texorpdfstring{\mr{बिघडणे} (bighaḍṇe)}{बिघडणे (bighaḍṇe)}]{\mr{बिघडणे} (bighaḍṇe) — to get worse, to break down}

\label{lesson:MR-C165-bighadne}



\begin{quote}
Before the new one: say the Marathi for to miss; to go wrong, then the Marathi for to improve.
\end{quote}

\subsection*{You'll want to know: \mr{बिघडणे}}
\textbf{\mr{बिघडणे}} — \emph{bighaḍṇe} — \textquotedblleft{}to get worse, to break down\textquotedblright{}.

\begin{grammarlens}[title={What you've built}]
One more word for this chapter.
\end{grammarlens}

\subsection*{Your turn}
\begin{itemize}
\raggedright
  \item \emph{Say it:} \emph{bighaḍṇe}
  \item \emph{Say it:} \emph{bighaḍṇe}, once more
  \item \emph{Say it:} say \emph{bighaḍṇe}
  \item \emph{Recall:} say the Marathi for to miss; to go wrong, then the Marathi for to improve, then say \emph{bighaḍṇe} again
\end{itemize}

\begin{tcolorbox}[breakable,colback=teal!4,colframe=teal!35!black,title={Before you move on}]
What does \mr{बिघडणे} mean? (\textquotedblleft{}To get worse, to break down\textquotedblright{}.) Say it once more.
\end{tcolorbox}
\section[\texorpdfstring{\mr{निरीक्षण} \mr{करणे} (nirīkṣaṇ karṇe)}{निरीक्षण करणे (nirīkṣaṇ karṇe)}]{\mr{निरीक्षण} \mr{करणे} (nirīkṣaṇ karṇe) — to observe}

\label{lesson:MR-C165-niriksan-karne}



\begin{quote}
Before the new one: say the Marathi for to improve, then the Marathi for to get worse.
\end{quote}

\subsection*{You'll want to know: \mr{निरीक्षण} \mr{करणे}}
\textbf{\mr{निरीक्षण} \mr{करणे}} — \emph{nirīkṣaṇ karṇe} — \textquotedblleft{}to observe\textquotedblright{}.

\begin{grammarlens}[title={What you've built}]
One more word for this chapter.
\end{grammarlens}

\subsection*{Your turn}
\begin{itemize}
\raggedright
  \item \emph{Say it:} \emph{nirīkṣaṇ karṇe}
  \item \emph{Say it:} \emph{nirīkṣaṇ karṇe}, once more
  \item \emph{Say it:} say \emph{nirīkṣaṇ karṇe}
  \item \emph{Recall:} say the Marathi for to improve, then the Marathi for to get worse, then say \emph{nirīkṣaṇ karṇe} again
\end{itemize}

\begin{tcolorbox}[breakable,colback=teal!4,colframe=teal!35!black,title={Before you move on}]
What does \mr{निरीक्षण} \mr{करणे} mean? (\textquotedblleft{}To observe\textquotedblright{}.) Say it once more.
\end{tcolorbox}
\section[\texorpdfstring{\mr{भाग} \mr{घेणे} (bhāg gheṇe)}{भाग घेणे (bhāg gheṇe)}]{\mr{भाग} \mr{घेणे} (bhāg gheṇe) — to take part}

\label{lesson:MR-C165-bhag-ghene}



\begin{quote}
Before the new one: say the Marathi for to get worse, then the Marathi for to observe.
\end{quote}

\subsection*{You'll want to know: \mr{भाग} \mr{घेणे}}
\textbf{\mr{भाग} \mr{घेणे}} — \emph{bhāg gheṇe} — \textquotedblleft{}to take part\textquotedblright{}.

\begin{grammarlens}[title={What you've built}]
One more word for this chapter.
\end{grammarlens}

\subsection*{Your turn}
\begin{itemize}
\raggedright
  \item \emph{Say it:} \emph{bhāg gheṇe}
  \item \emph{Say it:} \emph{bhāg gheṇe}, once more
  \item \emph{Say it:} say \emph{bhāg gheṇe}
  \item \emph{Recall:} say the Marathi for to get worse, then the Marathi for to observe, then say \emph{bhāg gheṇe} again
\end{itemize}

\begin{tcolorbox}[breakable,colback=teal!4,colframe=teal!35!black,title={Before you move on}]
What does \mr{भाग} \mr{घेणे} mean? (\textquotedblleft{}To take part\textquotedblright{}.) Say it once more.
\end{tcolorbox}
